\chapter{Trapeziectomy}

\section*{What Is This Surgery?}
Trapeziectomy is a surgical procedure that involves the complete excision of the trapezium bone, a key carpal bone located at the base of the thumb. This intervention is primarily indicated for patients suffering from severe pain and functional limitations due to advanced osteoarthritis of the first carpometacarpal (CMC) joint, commonly referred to as basal joint arthritis. The surgical removal of the trapezium alleviates the painful bone-on-bone contact that characterizes this degenerative condition. Following the excision, the resultant space may be filled with a tendon graft or reconstructed with a ligament to provide stability and prevent the thumb metacarpal from subsiding into the wrist. This procedure is recognized as a definitive treatment for end-stage degeneration of the thumb's basal joint, significantly improving the quality of life for affected individuals.

\section*{Alternative Names}
\begin{itemize}
  \item Trapezium Excision
  \item Basal Joint Arthroplasty
  \item CMC Arthroplasty
  \item Excisional Arthroplasty of the Thumb CMC Joint
  \item Ligament Reconstruction Tendon Interposition (LRTI) Arthroplasty
  \item Suspension Arthroplasty
\end{itemize}

\section*{Relevant Surgical Anatomy}
The trapezium is one of the eight carpal bones in the wrist and is located at the base of the first metacarpal bone. It articulates with the first metacarpal, scaphoid, and trapezoid bones, forming the CMC joint. The CMC joint is critical for thumb opposition and grip strength. 

The arterial supply to the trapezium is primarily from the radial artery, which provides branches that supply the surrounding soft tissues and bones. Venous drainage occurs via the dorsal venous network of the hand, which drains into the cephalic and basilic veins. 

Nerve supply is provided by the median nerve, which innervates the thenar muscles and provides sensation to the palmar aspect of the thumb. The radial sensory nerve also supplies sensation to the dorsal aspect of the thumb. 

Lymphatic drainage from the hand follows the venous system, with lymphatic vessels draining into the axillary lymph nodes. 

Key surgical landmarks include:

\begin{itemize}
  \item McBurney's Point: While not directly related, it is a reference point for anatomical orientation.
  \item Calot's Triangle: Important for understanding the vascular and nerve supply in adjacent structures.
\end{itemize}

Understanding these anatomical relationships is crucial for minimizing complications during surgery.

\section*{Surgical Principle \& Rationale}
The primary surgical principle of trapeziectomy is to remove the source of pain associated with advanced thumb carpometacarpal (CMC) osteoarthritis. In this condition, the degeneration of articular cartilage leads to painful bone-on-bone contact, resulting in inflammation and significant discomfort during thumb movement and pinch activities. 

By excising the trapezium, the painful joint surfaces are permanently separated, alleviating the symptoms associated with the condition. The challenge post-excision is to prevent the first metacarpal from migrating proximally into the space left by the trapezium, which could lead to thumb shortening and instability. 

To address this, techniques such as interposition arthroplasty or ligament reconstruction are employed. In interposition arthroplasty, a tendon graft, often from the flexor carpi radialis (FCR) or abductor pollicis longus (APL), is used to fill the void left by the trapezium, providing cushioning and stability. Ligament reconstruction involves using a tendon graft to recreate the stabilizing ligaments of the CMC joint, ensuring the metacarpal remains properly positioned. These adjunct procedures aim to restore stability, preserve thumb length, and facilitate the formation of a new, pain-free pseudoarthrosis.

\section*{History \& Surgical Evolution}
The surgical approach to treating painful joints has evolved significantly since the early 20th century. Initially, procedures such as arthrodesis (fusion) were performed to stabilize the thumb CMC joint, but this approach sacrificed mobility. Simple trapeziectomy was introduced but often resulted in thumb shortening and instability.

A pivotal advancement occurred in the 1970s with the introduction of ligament reconstruction tendon interposition (LRTI) arthroplasty, pioneered by surgeons G. Eaton and J. Littler. They refined the technique of trapeziectomy by excising the trapezium and using a portion of the flexor carpi radialis tendon to reconstruct the anterior oblique ligament and interpose tissue into the void. This combined approach significantly improved surgical outcomes by providing both pain relief and stability. Over the years, variations of this technique have emerged, including different tendon sources and the use of synthetic materials, but the core principle of trapeziectomy with interposition or suspension remains a cornerstone of surgical management for advanced thumb basal joint arthritis.

\section*{Clinical Indications}
Trapeziectomy is indicated in several clinical scenarios, including:

\begin{itemize}
  \item Severe, persistent pain at the base of the thumb that significantly interferes with daily activities.
  \item Failure of comprehensive non-operative management, including splinting, activity modification, non-steroidal anti-inflammatory drugs (NSAIDs), and corticosteroid injections.
  \item Progressive functional impairment, such as weakness of pinch and grip strength, difficulty with fine motor tasks, and inability to open jars or turn keys.
  \item Radiographic evidence of advanced carpometacarpal (CMC) joint osteoarthritis, typically Eaton-Littler Stage III or IV, characterized by significant joint space narrowing, subchondral sclerosis, osteophytes, and sometimes subluxation.
  \item Persistent tenderness and crepitus upon palpation and axial loading of the CMC joint (positive grind test).
  \item Progressive adduction contracture of the thumb metacarpal, limiting thumb abduction and opposition.
  \item Instability or subluxation of the CMC joint that contributes to pain and dysfunction.
\end{itemize}

These indications highlight the necessity for surgical intervention when conservative measures fail to provide relief.

\section*{Clinical Positioning}
Trapeziectomy is primarily considered a first-line surgical treatment for advanced, symptomatic osteoarthritis of the thumb carpometacarpal (CMC) joint. It is typically recommended after a thorough trial of conservative management has failed to yield adequate pain relief or functional improvement. 

Conservative measures, including activity modification, splinting, oral medications, and corticosteroid injections, are the first line of treatment for early to moderate stages of CMC arthritis. However, once the disease progresses to severe, end-stage degeneration with significant pain and functional loss, trapeziectomy, often combined with ligament reconstruction and/or tendon interposition, becomes the primary surgical option. While it can theoretically serve as a secondary or salvage procedure if other arthroplasty techniques (e.g., implant arthroplasty) have failed, its most common role is as the initial definitive surgical intervention for advanced basal joint arthritis.

\section*{Surgical Technique \& Procedure}
The trapeziectomy procedure is typically performed under regional anesthesia (e.g., axillary block) with sedation or general anesthesia, depending on patient preference and surgeon discretion. The patient is positioned supine with the arm extended on a hand table, and a tourniquet is applied to the upper arm to ensure a bloodless field.

The surgical steps include:

\begin{itemize}
  \item A curvilinear or straight incision is made over the dorsoradial aspect of the thumb carpometacarpal joint, carefully avoiding injury to branches of the radial sensory nerve.
  \item The underlying soft tissues are dissected to expose the trapezium bone and its articulations with the first metacarpal, scaphoid, and trapezoid.
  \item The capsule of the CMC joint is incised, and the trapezium is meticulously dissected free from its surrounding ligamentous attachments and excised using osteotomes, rongeurs, and burrs.
  \item Care is taken to remove the entire bone while preserving the surrounding structures.
  \item Once the trapezium is removed, the space is prepared for stabilization. In a common technique like Ligament Reconstruction Tendon Interposition (LRTI), a strip of the flexor carpi radialis (FCR) tendon is harvested from the forearm, leaving its distal attachment intact.
  \item This tendon strip is then passed through a drill hole in the base of the first metacarpal, woven through the remaining capsule and ligaments, and secured to itself or to the adjacent flexor carpi radialis tendon, creating a suspension sling.
  \item The remaining portion of the FCR tendon is rolled up and placed into the space created by the trapeziectomy, acting as an interpositional cushion.
  \item After ensuring adequate stability and hemostasis, the wound is irrigated, and the capsule and skin are closed in layers.
  \item A bulky dressing and a thumb spica splint or cast are applied to immobilize the thumb and wrist in a functional position for several weeks to allow for healing and stabilization.
\end{itemize}

This detailed approach ensures optimal outcomes and minimizes complications.

\section*{Preoperative Workup \& Preparation}
Prior to trapeziectomy, a comprehensive pre-operative evaluation is essential to ensure patient safety and optimize surgical outcomes.

\begin{itemize}
  \item \textbf{Medical Evaluation:} A thorough medical history and physical examination are performed, assessing for co-morbidities, allergies, and current medications.
  \item \textbf{Laboratory Tests:} Routine pre-operative blood tests, such as a complete blood count (CBC), basic metabolic panel, and coagulation studies, are typically ordered.
  \item \textbf{Imaging Studies:} Standard X-rays of the hand, including AP, lateral, and specific views like Robert's view or stress views, are crucial to confirm the diagnosis of CMC osteoarthritis, assess its severity (e.g., Eaton-Littler staging), and rule out other pathologies.
  \item \textbf{Anesthesia Consultation:} Patients undergo a pre-anesthesia evaluation to determine the most appropriate type of anesthesia and to assess any risks.
  \item \textbf{Medication Management:} Patients may need to adjust or temporarily discontinue certain medications, particularly anticoagulants or antiplatelet agents, several days before surgery to minimize bleeding risk.
  \item \textbf{NPO Status:} Patients are instructed to refrain from eating or drinking for a specified period (typically 6-8 hours) before surgery to prevent aspiration during anesthesia.
  \item \textbf{Prophylactic Antibiotics:} Intravenous antibiotics are administered shortly before the incision to reduce the risk of surgical site infection.
  \item \textbf{Informed Consent:} A detailed discussion with the surgeon covers the risks, benefits, alternatives, and expected recovery process of the procedure, and the patient provides informed consent.
\end{itemize}

This thorough preparation is critical for minimizing complications and ensuring a successful surgical outcome.

\section*{Contraindications \& Precautions}
\textbf{Absolute Contraindications:}
\begin{itemize}
  \item Active infection in the hand, wrist, or systemic infection, as this significantly increases the risk of surgical site infection and complications.
  \item Severe systemic illness or uncontrolled medical conditions (e.g., unstable cardiac disease, uncontrolled diabetes) that would make the patient an unacceptable anesthetic or surgical risk.
  \item Very mild or early-stage carpometacarpal (CMC) arthritis that has not failed a comprehensive trial of conservative management.
  \item Significant neurovascular compromise in the affected limb.
\end{itemize}

\textbf{Relative Contraindications and Precautions:}
\begin{itemize}
  \item Inflammatory arthritis (e.g., rheumatoid arthritis, psoriatic arthritis) may require different surgical considerations or techniques due to generalized joint laxity or bone quality issues.
  \item Significant generalized ligamentous laxity, which might compromise the stability of the reconstruction.
  \item Patient non-compliance with post-operative immobilization and rehabilitation protocols, which are critical for successful outcomes.
  \item Very high functional demands or heavy manual labor, where alternative procedures like CMC joint fusion might be considered for maximum strength and stability, albeit at the expense of motion.
  \item Severe osteoporosis, which could complicate tendon anchor placement or lead to intraoperative fractures.
  \item Concomitant wrist pathology that might influence the surgical approach or recovery.
\end{itemize}

Understanding these contraindications is essential for patient selection and optimizing surgical outcomes.

\section*{Complications \& Risks}
While trapeziectomy is generally a safe and effective procedure, it carries potential risks and complications that must be communicated to patients. These can be broadly categorized as general surgical risks and procedure-specific risks.

\begin{itemize}
  \item \textbf{General Surgical Risks:}
    \begin{itemize}
      \item \textbf{Bleeding:} Hematoma formation at the surgical site.
      \item \textbf{Infection:} Superficial wound infection or deep joint infection, requiring antibiotics or further surgery.
      \item \textbf{Anesthesia Complications:} Reactions to anesthetic agents, nausea, vomiting, respiratory issues, or nerve injury related to regional blocks.
      \item \textbf{Scarring:} Hypertrophic or keloid scarring, scar tenderness.
    \end{itemize}
  \item \textbf{Procedure-Specific Risks:}
    \begin{itemize}
      \item \textbf{Nerve Injury:} Damage to branches of the radial sensory nerve, leading to numbness, tingling, or painful neuroma formation on the back of the thumb or hand.
      \item \textbf{Thumb Shortening (Subsidence):} Proximal migration of the first metacarpal into the space left by the excised trapezium, potentially leading to decreased pinch strength and cosmetic deformity, especially if interposition/suspension is inadequate.
      \item \textbf{Persistent Pain:} Inadequate pain relief or development of new pain due to incomplete excision, instability, or other factors.
      \item \textbf{Stiffness or Limited Range of Motion:} Restricted movement of the thumb or wrist, despite rehabilitation.
      \item \textbf{Weakness:} Decreased pinch or grip strength, which may persist despite pain relief.
      \item \textbf{Complex Regional Pain Syndrome (CRPS):} A rare but severe chronic pain condition affecting the limb, characterized by pain, swelling, stiffness, and skin changes.
      \item \textbf{Tendon-Related Complications:} If a tendon graft is used, risks include tendon rupture, tendinitis, or donor site morbidity.
      \item \textbf{Failure of Reconstruction:} The interposition or suspension technique may fail to adequately stabilize the thumb, leading to instability or persistent symptoms.
      \item \textbf{Non-union/Malunion:} Not applicable to trapeziectomy itself, but relevant if fusion was considered or if an implant was used and failed.
    \end{itemize}
\end{itemize}

Awareness of these complications is vital for informed consent and patient education.

\section*{Postoperative Recovery Timeline}
Following trapeziectomy, the patient is closely monitored in the post-anesthesia care unit (PACU) for any immediate complications. Pain management protocols are initiated, and the hand is typically immobilized in a bulky dressing and a thumb spica splint or cast, which extends from the forearm to the thumb, holding it in a position of rest and slight abduction.

In the first 24-48 hours, patients are usually discharged home with detailed instructions for pain management, wound care, and elevation of the hand to minimize swelling. The initial splint or cast is generally maintained for approximately 4 to 6 weeks to allow for soft tissue healing and stabilization of the thumb metacarpal. During this period, patients are advised to keep the dressing dry and clean and to avoid any strenuous activities or movements of the thumb.

After the initial immobilization phase, the splint is removed, and a structured hand therapy program is initiated. This program is crucial for gradually restoring range of motion, improving strength, and regaining functional use of the thumb. Full recovery, including significant pain relief and return to most daily activities, can take anywhere from 3 to 6 months, with continued improvement possible for up to a year.

\section*{Surgical Findings \& Pathology}
During the procedure, the surgeon documents the condition of the trapezium and surrounding structures. The excised trapezium is typically sent to pathology, where microscopic examination confirms the presence of degenerative changes characteristic of osteoarthritis, such as cartilage erosion, subchondral bone sclerosis, and osteophyte formation. This pathological confirmation validates the pre-operative diagnosis and provides insight into the severity of the disease.

The surgical findings may also include the assessment of the surrounding ligaments and tendons, which can influence the choice of reconstruction techniques used during the procedure.

\section*{Expected Postoperative Course}
A normal postoperative course following trapeziectomy is characterized by a significant reduction or elimination of pain at the base of the thumb, particularly during pinch and grip activities. Patients should experience improved functional use of their hand, including enhanced pinch strength, better dexterity, and an increased ability to perform daily tasks without discomfort. 

The stability of the thumb carpometacarpal joint should be maintained, and the thumb metacarpal should not have significantly subsided into the wrist. While some degree of thumb shortening is common, it should not be functionally limiting. Ultimately, the expected outcome is a more comfortable and functional thumb, allowing the patient to return to activities that were previously painful or impossible.

\section*{Postoperative Care \& Follow-up}
Post-operative follow-up after trapeziectomy is crucial for monitoring healing, guiding rehabilitation, and addressing any potential complications.

\begin{itemize}
  \item \textbf{Initial Post-operative Visits:} Typically, the first follow-up occurs within 1-2 weeks for wound check and possibly splint adjustment or cast change. Suture removal usually takes place around 2-3 weeks post-surgery.
  \item \textbf{Immobilization and Splinting:} The thumb spica splint or cast is generally maintained for 4-6 weeks. After this period, a removable splint may be used for protection during activities.
  \item \textbf{Hand Therapy/Rehabilitation:} This is a critical component of recovery, usually starting after the initial immobilization period. A certified hand therapist will guide exercises to restore range of motion, improve strength, and enhance functional use of the thumb. Therapy may continue for several months.
  \item \textbf{Activity Resumption:} Patients are gradually cleared to resume light activities, progressing to more strenuous tasks under the guidance of their surgeon and therapist. Heavy lifting and forceful gripping are typically restricted for several months.
  \item \textbf{Imaging:} Routine post-operative X-rays are usually not required unless there are concerns about instability, persistent pain, or other complications.
  \item \textbf{Warning Signs:} Patients are educated on warning signs to report to their doctor immediately, such as increasing pain, redness, swelling, fever, drainage from the wound, or new numbness/tingling.
\end{itemize}

Regular follow-up appointments with the surgeon continue for several months to a year to monitor progress and ensure optimal long-term outcomes.

\section*{Epidemiology}
Osteoarthritis of the thumb carpometacarpal (CMC) joint, for which trapeziectomy is a primary surgical treatment, is a highly prevalent condition, particularly affecting older adults. It is estimated to affect approximately 1 in 3 post-menopausal women and 1 in 8 men over the age of 50. The incidence significantly increases with age, with radiographic evidence of CMC arthritis found in up to 80\% of individuals over 75 years old. 

Women are disproportionately affected, with a female-to-male ratio ranging from 6:1 to 10:1, suggesting hormonal factors may play a role in its pathogenesis. While many individuals with radiographic evidence of CMC arthritis remain asymptomatic, a substantial portion develops debilitating pain and functional impairment requiring intervention. Trapeziectomy is one of the most commonly performed hand surgeries for arthritis, reflecting the high frequency of symptomatic basal joint disease. Mortality directly related to trapeziectomy is exceedingly rare, but morbidity from persistent pain, stiffness, or nerve injury, while low, can significantly impact quality of life.

\section*{Outcomes \& Prognosis}
The outcomes following trapeziectomy, particularly when combined with ligament reconstruction and tendon interposition (LRTI), are generally excellent, with high rates of patient satisfaction. Studies consistently report good to excellent results in 85\% to 95\% of patients, characterized by significant pain relief and improved functional use of the thumb. Patients typically experience a substantial reduction in pain during pinch and grip activities, leading to an enhanced ability to perform daily tasks. 

While some degree of thumb shortening is common, it is rarely functionally significant. The prognosis for long-term pain relief and functional improvement is favorable, with results often sustained for many years. Recurrence, in the sense of the trapezium growing back, does not occur. However, potential long-term issues can include persistent mild pain, some degree of residual weakness, or stiffness. Factors influencing prognosis include the patient's pre-operative severity of arthritis, adherence to post-operative rehabilitation, the presence of other hand pathologies, and overall health status. While complete restoration of pre-arthritic strength may not always be achieved, the primary goal of a stable, pain-free, and functional thumb is typically met, significantly improving the patient's quality of life.

\section*{References}
\begin{enumerate}
  \item MedlinePlus. Trapeziectomy. U.S. National Library of Medicine; 2025.
  \item Green DP, Hotchkiss RN, Pederson WC, Wolfe SW. Green's Operative Hand Surgery. 7th ed. Elsevier; 2017.
  \item American Academy of Orthopaedic Surgeons (AAOS). Management of Osteoarthritis of the Thumb Carpometacarpal Joint: Evidence-Based Clinical Practice Guideline. AAOS; 2014.
  \item Canale ST, Beaty JH. Campbell's Operative Orthopaedics. 14th ed. Elsevier; 2021.
\end{enumerate}

\clearpage